   \documentclass{article}
    
    \textwidth=6.5in
    \textheight=8.6in
    \voffset=-54pt
    \oddsidemargin=0in
    \evensidemargin=0in
    \setlength{\parskip}{8pt}
    \setlength{\parindent}{0pt}
    
  
    \usepackage[normalem]{ulem}
    \usepackage{microtype}
    
    
    \usepackage{scalerel,amssymb}
\def\mcirc{\mathbin{\scalerel*{\circ}{j}}}
\def\msquare{\mathord{\scalerel*{\Box}{gX}}}

    \usepackage{diffcoeff}

    \usepackage{amsmath, amssymb}
    \usepackage{ifthen}
    
    \usepackage[inline]{enumitem}
    \setlist{topsep=0pt}
    
    \usepackage{xcolor}

    \newcommand{\eqdots}{\setbox0=\hbox{=}%
       \mathrel{\hbox to\wd0{\hss\vdots\hss}}}
    \usepackage[percent]{overpic}
    \usepackage{pifont}
   \usepackage{siunitx} 
    
    % change hyperref options
    \PassOptionsToPackage{hyphens}{url}

\usepackage[colorlinks,linkcolor=blue,urlcolor=blue,breaklinks]{hyperref}

    
    \usepackage{graphicx}
    \usepackage{multicol}
    \usepackage{framed}
    \usepackage{array}
    \usepackage{soul}
    \usepackage{wrapfig}
    
    \usepackage{pgfplots}
    \pgfplotscreateplotcyclelist{mycolor}{%
      black, blue, red, green!70!black,magenta,
      purple, cyan, orange
    }
    \pgfplotsset{
      cycle list name=mycolor,
      axis lines=middle,
      every axis plot/.style={no marks, thick},
      every axis/.style={
        enlarge x limits=false,
        enlarge y limits=auto,
        xlabel style={at={(current axis.right of origin)},anchor=west},
        ylabel style={at={(current axis.above origin)},anchor=south},
        % extra x and y ticks for asymptotes
        extra x tick style={grid=major, grid style={dashed,black}},
        extra y tick style={grid=major, grid style={dashed,black}},
        /pgf/number format/1000 sep={},
      },
      tick label style = {font=\footnotesize, fill=white},
      scaled ticks=false,
      scaled y ticks=false,
      unbounded coords=jump,
      samples=101,
      scatter/classes={
        open={mark=*,draw=black,fill=white, mark size=2, thin},
        closed={mark=*,draw=black,fill=black, mark size=2, thin}
      },
    }
      
    
    \usepackage{tikz}
    \usetikzlibrary{
      % enables the snake paths
      decorations.pathmorphing,%
      % enables bigger arrows
      decorations.markings,%
      decorations.pathreplacing,%
      decorations.text,%
      calc,%
      patterns, patterns.meta,
      shapes.geometric,%
      arrows,
      decorations.shapes,
      positioning,
      plotmarks, angles, quotes,
      pgfplots.groupplots
    }
    \tikzset{>=latex} % sets default arrow type in tikz
    \tikzset{font=\small}
    
      
    \usepackage{multirow, bigdelim}
    
    % change to \usepackage[sols]{handout} to see solutions
    \usepackage[]{handout}
        
    \newcommand{\lheq}{\stackrel{\text{L'H}}{=}}
\newcommand{\nolheq}{\mathrel{\makebox[\widthof{$\lheq$}]{=}}}

    \definecolor{skyblue}{rgb}{0, 0.5, 1}
    \usepackage{tcolorbox}
\newenvironment{feedback}{%
  \begin{tcolorbox}[colframe=skyblue, colback=skyblue!10!white]
  }{\end{tcolorbox}}
      
    
    \newcommand\iscurrentenumi[1]{%
      \ifthenelse{\equal{#1}{\theenumi}}{}{#1}}
    \setenumerate[1]{ref=\#\arabic*}
    \setenumerate[2]{ref=\protect\iscurrentenumi{\theenumi}(\alph*)}

    
      
    
    \newlength{\tipmargin}
    \makeatletter
    
    \makeatother
    
      
    \ifthenelse{\isundefined{\C}}{%
      \newcommand{\C}{\mathbb{C}}%
    }{\renewcommand{\C}{\mathbb{C}}}
    \ifthenelse{\isundefined{\R}}{%
      \newcommand{\R}{\mathbb{R}}%
    }{\renewcommand{\R}{\mathbb{R}}}
    
    \definecolor{purple}{rgb}{0.5,0,1.0}
    \pgfplotsset{holdot/.style={color=red,fill=white,only marks,mark=*}}
    \pgfplotsset{soldot/.style={color=black,fill=black,only marks,mark=*}}
\pgfplotsset{compat=1.13}
\usetikzlibrary{quotes,arrows.meta}

    \pgfplotsset{holdot/.style={color=red,fill=white,only marks,mark=*}}

\newcommand{\answerbox}[3][]{
    \fbox{\begin{minipage}[t][#2][c]{#3}
        Final answer: #1
    \end{minipage}}
    }

\newcommand{\answerboxd}[3][]{
    \fbox{\begin{minipage}[t][#2][c]{#3}
        #1
    \end{minipage}}
    }
    
\begin{document}

 

 \noindent
{\large \mbox{} \hfill \bf Name:\underline{\hspace{2.8in}}}

 {\mbox{} \hfill \it (Please write your name neatly in print.)}


 
\medskip
 
\noindent
{\large\bf
Trig Skills Check \\ Math Mb Spring 2025}
\noindent






\iffalse
\begin{center}
\begin{tabular}{|c|c|c|}
\hline
Problem & Points & Score\\ \hline
& &\\
1 &  & \\ 
\hline
&& \\
2&  &  \\ 
\hline
& &\\
3&&  \\ 
 \hline
& &\\
Total&  &  \\ 
& &\\ \hline

\end{tabular}
\end{center}
\fi


%{\large \em Please circle your section:}


%\begin{center}
%{\small
%\begin{tabular}{ccccccccc}
%Erica Dinkins & Matt Cavallo &  Shanelle Davis  & Robin Gottlieb & 	Mike Koske & Hakim Walker \\
%MWF 9am  & MWF 9am &  MWF 10:30am  &MWF 10:30am & MWF 12pm   & MWF 12pm\\
%\\\\
% &&  &     \\
% &&  &
%\end{tabular}

%}
%\vspace{ -.5in} 
%{\small
%\begin{tabular}{cccccccc}
 %Hannah Constantin & Matt Cavallo & 	Justin Hancock  & Grant Barkley \\
 %MWF 12pm &  MWF 12pm  & MWF 1:30pm  & MWF 3pm  \\
%\\\\
 %&&       \\
 %&&  
%\end{tabular}


%}
%\end{center}


{\Large
\begin{center}\textbf{Directions---Please read carefully!} \end{center}}

You have 60 minutes to complete this Skills Check. \textbf{Calculators or any other aids are not permitted.}  


For each problem, write your final answer in the answer box provided. You may do work on any page and may ask for scratch paper, but \textbf{no work or writing outside of the answer box will be graded.}

Write neatly, as answers that cannot be read will not receive credit. 

   

\begin{center}\textbf{Good luck!}\end{center}


\newpage

\begin{enumerate}

\item Evaluate each expression.

    \begin{multicols}{2}
    \begin{enumerate}[label=(\alph*)]
        \item $\sin(3\pi/2)$
        \item $\cos(3\pi/4)$
    \end{enumerate}
    \end{multicols}

    \vfill 

        \begin{multicols}{2}
                \answerbox[\solonly{$-1$}]{0.8in}{2.8in}
                \answerbox[\solonly{$-\frac{\sqrt{2}}{2}$}]{0.8in}{2.8in}
        \end{multicols}

        \begin{multicols}{2}
    \begin{enumerate}[start=3,label=(\alph*)]
        \item $\cos(11\pi/3)$
        \item $\sin(4\pi/3)$
    \end{enumerate}
    \end{multicols}

    \vfill

        \begin{multicols}{2}
                \answerbox[\solonly{$\frac{1}{2}$}]{0.8in}{2.8in}
                \answerbox[\solonly{$-\frac{\sqrt{3}}{2}$}]{0.8in}{2.8in}
        \end{multicols}

    \begin{multicols}{2}
    \begin{enumerate}[start=5,label=(\alph*)]
        \item $\sin(13\pi/4)$
        \item $\cos(-11\pi/6)$
    \end{enumerate}
    \end{multicols}

    \vfill 

        \begin{multicols}{2}
                \answerbox[\solonly{$\frac{-\sqrt{2}}{2}$}]{0.8in}{2.8in}
                \answerbox[\solonly{$\frac{\sqrt{3}}{2}$}]{0.8in}{2.8in}
        \end{multicols}



\newpage 

\item Solve the equation $\sin(\theta)=-\frac{1}{2}$ on the interval $[0,2\pi)$.

\vfill

\answerbox[\solonly{$\theta=\dfrac{7\pi}{6}, \dfrac{11\pi}{6}$}]{0.8in}{5.5in}

 \item $A$ is an angle such that $\sin A=-\frac{5}{6}$ and $\cos A < 0$. Find $\cos A$.

        \vfill

        \answerbox[\solonly{$\cos A = -\dfrac{\sqrt{11}}{6}$}]{0.8in}{5.5in}

\item Convert 6 degrees to radians. Simplify your answer fully.

\vfill

\answerbox[\solonly{$\dfrac{\pi}{30}$ radians}]{0.8in}{5.5in}


\newpage 

\item
Consider the graph of the sinusoidal function $y=f(t)$ shown below.

\begin{center}
            \begin{tikzpicture}
          \begin{axis}[xtick=\empty, ytick=\empty, clip=false, width=4in, height=2in,
          xlabel=$t$, ylabel=$y$, ylabel style={rotate=-90}]
            \addplot+[domain=-2.25:2.25, samples=100]{-4*cos(deg(2*pi*x))-1};
            \addplot[soldot] coordinates {(0,-5) (1.5,3)};
            \node[left] at (0,-5) {$(0,-5)$};
            \node[above] at (1.5,3) {$(\frac{3}{2},3)$};
          \end{axis}
        \end{tikzpicture}
\end{center}
    \begin{enumerate}
        \item Find the midline, amplitude, and period.

        \vfill

\answerboxd[Midline: $y=\solonly{-1}$]{.5 in}{5.5 in}

\vspace{.1in}
    
\answerboxd[Amplitude $=\solonly{3}$]{.5 in}{5.5 in}

\vspace{.1in}

\answerboxd[Period $=\solonly{2}$]{.5 in}{5.5 in}

\vspace{.1in}

    \item Write a formula for the function.

    \vfill

    \answerbox[$f(t)=\solonly{-3\cos(\pi t)-1}$]{0.8in}{5.5in}
   
    \end{enumerate}
    
\newpage

\item Determine whether each statement is an identity (true for all values). Write ``True" or ``False."


    \begin{multicols}{2}
            \begin{enumerate}
        \item  $\sin(\theta)=-\sin(-\theta)$

        \item $\sin(AB)=\sin(A)\sin(B)$

    \end{enumerate}
    \end{multicols}

    \vfill 

        \begin{multicols}{2}
                \answerbox[\solonly{True}]{0.8in}{2.5in}
                
                \answerbox[\solonly{False}]{0.8in}{2.5in}
        
    \end{multicols}

    \begin{multicols}{2}
            \begin{enumerate}[start=3]
        \item $\sin^2\theta-1 = \cos^2\theta$
        \item $\sin(\theta+\frac{\pi}{2})=\cos(\theta)$
    \end{enumerate}
    \end{multicols}
\vfill 

    \begin{multicols}{2}
                \answerbox[\solonly{False}]{0.8in}{2.5in}
                
                \answerbox[\solonly{True}]{0.8in}{2.5in}
        
    \end{multicols}


\newpage 
        
\item Evaluate each expression. If the expression is undefined, write ``undefined."

    \begin{multicols}{2}
    \begin{enumerate}[label=(\alph*)]
        \item $\tan(7\pi/4)$
        \item $\tan(-7\pi/6)$
    \end{enumerate}
    \end{multicols}

    \vfill 

        \begin{multicols}{2}
                \answerbox[\solonly{$-1$}]{0.8in}{2.8in}
                \answerbox[\solonly{$-\dfrac{\sqrt{3}}{3}$}]{0.8in}{2.8in}
        \end{multicols}

        \begin{multicols}{2}
    \begin{enumerate}[start=3,label=(\alph*)]
        \item $\sec(5\pi/3)$
        \item $\csc(5\pi/2)$
    \end{enumerate}
    \end{multicols}

    \vfill

        \begin{multicols}{2}
                \answerbox[\solonly{$2$}]{0.8in}{2.8in}
                \answerbox[\solonly{$1$}]{0.8in}{2.8in}
        \end{multicols}

    \begin{multicols}{2}
    \begin{enumerate}[start=5,label=(\alph*)]
        \item $\cot(\pi/6)$
        \item $\csc(3\pi)$
    \end{enumerate}
    \end{multicols}

    \vfill 

        \begin{multicols}{2}
                \answerbox[\solonly{$\sqrt{3}$}]{0.8in}{2.8in}
                \answerbox[\solonly{undefined}]{0.8in}{2.8in}
        \end{multicols}

\newpage

\item The shadow of a vertical tower is 40 m long when the angle of elevation of the sun is \ang{34.6}. Find the height of the tower. You may leave unevaluated trigonometric expressions in your answer. Include units.

\begin{tikzpicture}

    % Define coordinates
    \coordinate (TowerBase) at (4, 0);
    \coordinate (TowerTop) at (4, 2); % Height of tower set to 4 units
    \coordinate (ShadowStart) at (0, 0); % Shadow starts from the origin
    \coordinate (Sun) at (6,3);

    \draw[fill=gray, draw=black] (4, 0) rectangle (4.3, 2);

    \draw[dashed] (Sun) -- (ShadowStart);
    \draw (-1,0) -- (6,0);
    \draw[ultra thick] (0,0) -- (4,0);
        
    % Draw angle arc
    \draw (0.5,0) arc[start angle=0,end angle=atan(0.5),radius=0.5];

    % Sun illustration
    \foreach \angle in {0,45,...,315}
        \draw[thick,orange] (6,3) -- ++(\angle:0.5);
    \draw[fill=yellow!80,draw=orange,thick] (6,3) circle (0.3);
    \node at (12:1.2) {\ang{34.6}};


\end{tikzpicture}

\vfill

\answerbox[\solonly{$(40.6\text{ m})\tan(\ang{34.6})$}]{0.8in}{5.5in}


\item A ladder reaches 60 feet up the side of a building when the angle between the ladder and the ground is \ang{50}. How long is the ladder? You may leave unevaluated trigonometric expressions in your answer. Include units.

\vspace{\stretch{2}}

\answerbox[\solonly{$(60\text{ ft})\csc(\ang{50})$}]{0.8in}{5.5in}

% \item A passenger in an airplane at an altitude of 10 kilometers looks out the window and sees a town to the east of the plane. The angle of depression to the town is \ang{22}. How many kilometers east of the plane is the town? You may leave unevaluated trigonometric expressions in your answer.

        
\end{enumerate}



\end{document}